The paper chooses a ``suitable subgraph'' of the current candidacy graph at each induction step, but states its invariant using the whole updated candidacy graph and suppresses the quantitative choice when future neighbors lie in different error cohorts. It also says that the two one-vertex size discrepancies have the same proof without recording the altered normalization or the deletion of the artificial row. The affected nodes are QuantitativeSuperRegularThinning, EqualSizeParameterTable, EqualSizeInductionInvariant, VariableSizeInductionBase, VariableSizeInductionStep, VariableSizeInductionIteration, EqualSizeTerminalConversion, NearEqualRegularTriplet, NearEqualRegularPairGoodCodegrees, NearEqualDegreeCodegreeRegularityCriterion, NearEqualFourGraphs, MatchingInducedLeftRowDeletion, and SpreadBlowupAuxiliary.

The formalization makes that suitable-subgraph choice persistent. For every unembedded cluster and every supported partial embedding it stores a spanning graph $B_f^j$ contained edgewise in the true updated candidacy graph $A_f^j$, with level-dependent parameters $(\beta_x,D_x)$. Before a transition it thins the current witness, each future witness pulled back through the source matching, and each relevant host pair. This permits one simultaneous good-edge graph whose losses are bounded by the finite sums over all future cohorts. One common perfect matching is then tested separately by the scalar Four-Graphs lemma. A successful matching-induced graph is transported into, and stored as a subgraph of, the true next candidacy graph; a nonneighbor retains its old witness.

This changes no compatibility condition and admits no new embedding: every selected matching edge lies in a witness, every witness lies in the true candidacy graph, and every transported output edge is proved to lie in the true updated candidacy graph. The parameter table depends only on $d,\Delta$ and finitely many Four-Graphs thresholds, so it is chosen before $r,n$ and the instance, exactly as required for the paper's constants. At time zero the original candidacy graphs are the witnesses. At time $r$ there is no unembedded cluster, hence every witness clause is vacuous. EqualSizeTerminalConversion then forgets all witness data and returns precisely the paper's distribution on compatible embeddings with the same $C/n$ vertex-spread conclusion. Thus the departure is proof bookkeeping that rigorously implements the paper's suitable-subgraph sentence and returns to its exact terminal claim.

For the unequal-size sentence, NearEqualRegularTriplet records the paper's actual common-neighborhood window using $|W_3|$, while NearEqualFourGraphs repeats the same matching argument on the equal pair $W_1,W_2$ and absorbs the sole ratios $|W_3|/m\in[1-1/m,1+1/m]$ into a preselected $O(1/m)$ slack. NearEqualDegreeCodegreeRegularityCriterion makes the corresponding final regularity calculation explicit. In the current-extra case, MatchingInducedLeftRowDeletion removes the artificial row and returns the literal residual parts, with unchanged density and only the reserved error enlargement; hence that row is never transported into the true candidacy graph. In the future-extra case the genuine unmatched row is retained and treated by the paper's prescribed trimming. These nodes alter only intermediate auxiliary graphs and constants. After deletion and transport, the induction invariant is again stated on the true source and host clusters, and the same terminal conversion forgets all witnesses and auxiliary rows. Thus this second bookkeeping departure rigorously realizes the paper's assertion that the rest of the proof is the same and returns independently to the exact paper conclusion.
